\originalchapter{离子阱、边带与量子操作}
\originalsection{射频囚禁与运动模式}
静电势满足Laplace方程, 不能在自由空间把同号电荷稳定困在三维势能极小值. Paul阱改用射频四极场. 一维方程变为Mathieu形式$\dd^2x/\dd\tau^2+(a_x-2q_x\cos2\tau)x=0$, $\tau=\Omega_{\rm rf}t/2$. 小$|q_x|$的稳定区有慢的世俗运动和快速微运动, 世俗频率$\omega_x\simeq\Omega_{\rm rf}\sqrt{a_x+q_x^2/2}/2$. 这个有效势只在稳定参数与时间尺度分离时成立.

多个离子因Coulomb排斥形成平衡阵列. 在平衡点展开势能二阶, Hessian矩阵的本征向量为集体正常模; 激光作用到一个离子也可耦合整个链的模式. 模频、参与因子与加热率决定门操作能否把运动自由度用作总线.

\originalsection{Lamb--Dicke展开与边带冷却}
离子位移$x=x_0(a+a^\dagger)$,$x_0=\sqrt{\hbar/(2m\omega_t)}$. 激光耦合含$\exp(\ii kx)$, 定义$\eta=kx_0$. 当$\eta\sqrt{2n+1}\ll1$, 展开为$1+\ii\eta(a+a^\dagger)$. 在激光失谐分别为$0,-\omega_t,+\omega_t$时, 保留慢项得载波$\sigma_+$、红边带$\sigma_+a$、蓝边带$\sigma_+a^\dagger$耦合, Rabi频率分别为$\Omega,\eta\Omega\sqrt n,\eta\Omega\sqrt{n+1}$.

红边带激发消耗一个声子, 再经近似不改运动量子数的光抽运回到内部基态, 构成冷却循环. $n=0$对红边带为暗态. 要分辨边带需$\Gamma_{\rm eff}\ll\omega_t$, 非共振加热与反冲给有限$\bar n$. 用率方程$\dot{\bar n}=-(A_--A_+)\bar n+A_+$得稳态$\bar n=A_+/(A_--A_+)$, 只有$A_->A_+$才冷却.

\originalsection{自旋依赖力与纠缠门}
同时驱动红蓝边带可产生$H_I=\hbar\chi S_x(a\e^{-\ii\delta t}+a^\dagger\e^{\ii\delta t})$. 由于不同时间的对易子正比于$S_x^2$并与$H_I$对易, Magnus展开在二阶结束. 演化含依赖自旋的位移$\alpha(t)=\chi(1-\e^{\ii\delta t})/\delta$和几何相位$\Phi(t)S_x^2$, $\Phi=\chi^2(\delta t-\sin\delta t)/\delta^2$, 相位符号依$H_I$约定. 在$t_g=2\pi/|\delta|$时$\alpha=0$, 运动轨迹闭合而自旋仍获得相互作用相位.

两离子取$S_x=(\sigma_x^{(1)}+\sigma_x^{(2)})/2$, 则$S_x^2=I/2+\sigma_x^{(1)}\sigma_x^{(2)}/2$. 选相位便可产生纠缠门. 闭合轨迹解释为何该门能减轻初始热运动的影响, 但模式漂移、加热、不同模式参与与非Lamb--Dicke修正仍会使末态与运动纠缠. 完整算法和纠错见量子信息卷.
